
\ifintegriereaufgaben
\section*{Partialbruchzerlegung Aufgaben}\label{MA2-partialbruchzerlegung-aufgaben}

Regeln und Beispiele, die zur Bearbeitung der Aufgaben hilfreich sind, finden Sie im Skript\wennreferenzvorhanden{MA2-partialbruchzerlegung}{ (Abschnitt \abschnittsverweis{MA2-partialbruchzerlegung})} und, falls Sie nicht weiterkommen, schauen Sie \hyperref[MA2-partialbruchzerlegung-loesungen]{hier}.
\else
\section{Partialbruchzerlegung Aufgaben}\label{MA2-partialbruchzerlegung-aufgaben}

Regeln und Beispiele zu diesen Aufgaben finden Sie im Skript. Die L�sungen finden Sie \hyperref[MA2-partialbruchzerlegung-loesungen]{hier}.
\fi

\begin{exercisebox}[Faktorisierung]
	Zerlegen Sie die folgenden Polynome weitestgehend in {Linearfaktoren}.
	\begin{eqnarray*}
			(a)\;4y^3-8y^2-3y+9 \quad (b)\; a^4+3a^2-10 \quad (c) \; x^4+2x^3-35x^2  \quad (d)\; -64+b^2
	\end{eqnarray*}
\end{exercisebox}

\begin{exercisebox}[Polynomdivision]
	F�hren Sie folgende {Polynomdivisionen} aus.
	\begin{eqnarray*}
		\begin{array}{l l }
			(a)\;(x^2+x):(x-1) \quad & (b)\; (x^3-x^2-x+1):(x-1) \\[2ex] (c) \; (x^3+x^2-5x+3) : (x-1)^2  \quad & (d)\; (x^4+x^3-2x) :(x-1)
		\end{array}
	\end{eqnarray*}
\end{exercisebox}
\begin{exercisebox}[Gau�-Algorithmus]
	L�sen Sie die folgenden Gleichungsysteme:
	\begin{eqnarray*}
		(a)\; \left\{ \begin{array}{cccccrcl}
	 &A_1 &  &      & + & B_2 & =& 0 \\
	 &4A_1& +& 4A_2 & + & 2B_2& =& -2 \\
	 &4A_1&- & 2A_2 & + & 8B_2& =& -38
 \end{array} \right.\quad \quad
		(b)\;  \left\{ \begin{array}{crcrcrcl}
	 -&x & + & y & +& z &=& 0 \\
	 &x & -&3y & -&2z & =& 5 \\
	 &5x & +&y & +&z &=& -6
 \end{array} \right.
 \end{eqnarray*}
\end{exercisebox}

\begin{exercisebox}[Partialbruchzerlegung]
	F�hren Sie eine Partialbruchzerlegung durch.
	\begin{eqnarray*}
		(a)\; \frac{2x^3-x^2-6x+2}{(x-5)^2(x^2+3x+4)} \quad\quad (b)\; \displaystyle{ \frac{x+1}{x(x^2+1)^2}}
\quad & (d)\; \displaystyle{\frac{x^2-5x}{(x-2)^3}}
		\end{eqnarray*}
\end{exercisebox}


\begin{exercisebox}[Partialbruchzerlegung]
	F�hren Sie eine Partialbruchzerlegung durch. Hinweis: Die Nennerpolynome haben mindestens eine ganzzahlige Nullstellen.
	\begin{eqnarray*}
		\begin{array}{l l l }
			(a) \; \displaystyle{ \frac{1}{x^3+x}}  \quad\quad & (b)\; \displaystyle{ \frac{(4x^5-6x+13)(x-1)(x^2-1)}{x^3-x^2-x+1}} (c) \;  \displaystyle{ \frac{2x^2+x+1}{x^3+x^2-5x+3}} \\[2ex]
			(d)\; \displaystyle{ \frac{-x^4-2x^3+7x^2-5x+3}{ x^2+3x-4}} \quad\quad & (e)\; \displaystyle{ \frac{x^2+x+1}{x^3-3x+2}}
		\end{array}
	\end{eqnarray*}
\end{exercisebox}

\begin{exercisebox}[Partialbruchzerlegung] % Papula Band 1, Abschnitt V.8, Aufgabe 7
	F�hren Sie eine Partialbruchzerlegung durch.
	\begin{eqnarray*}
		\begin{array}{l l l}
			(a)\; \displaystyle{ \frac{1}{ x^2-a^2}}, \; a>0 \;  \quad \quad & (b)\; \displaystyle{ \frac{4x^3}{x^3+2x^2-x-2}} \quad\quad & (c) \; \displaystyle{ \frac{3z}{z^3-3z^2-4}} \\[2ex]
			(d)\; \displaystyle{\frac{4x-2}{x^2-2x-63}} \; & (e)\; \displaystyle{\frac{2x+1}{x^3-6x^2+9x} }  &
		\end{array}
	\end{eqnarray*}
\end{exercisebox}

\begin{exercisebox}[Anwendung Partialbruchzerlegung]
	Berechnen Sie folgende Integrale. Hinweis: Nutzen Sie die Ergebnsise aus der vorigen Aufgabe.
	\begin{eqnarray*}
		\begin{array}{l l l}
			(a)\; \displaystyle{ \int \frac{-x^4-2x^3+7x^2-5x+3}{ x^2+3x-4}} \, \mathrm{d}x \quad & (b)\; \displaystyle{\int \frac{x^2+x+1}{x^3-3x+2}}\,\mathrm{d}x  & (c) \; \displaystyle{\int \frac{1}{x^3+x}}\,\mathrm{d}x \\[1ex]
			(d)\; \displaystyle{\int \frac{x^2-5x}{(x-2)^3}} \,\mathrm{d}x &(e)\; \displaystyle{ \int \frac{1}{ x^2-a^2}} \, \mathrm{d}x, \; a>0 & (f)\; \displaystyle{\int \frac{4x^3}{x^3+2x^2-x-2}}\,\mathrm{d}x \\[2ex]
			(g) \; \displaystyle{\int \frac{3z}{z^3-3z^2-4}}\,\mathrm{d}x  & (h)\; \displaystyle{\int \frac{4x-2}{x^2-2x-63}}\, \mathrm{d}x \quad & (i)\; \displaystyle{\int \frac{2x+1}{x^3-6x^2+9x}}  \, \mathrm{d}x
		\end{array}
	\end{eqnarray*}
\end{exercisebox}

